% COMS 6998 Networks for AI -- Final project proposal template.
% Aim for about 1-2 pages. The headings are a guide, not a strict form:
% adapt them to your project. Submit the compiled PDF on Canvas.
\documentclass[11pt]{article}
\usepackage[margin=1in]{geometry}
\usepackage[T1]{fontenc}
\usepackage{lmodern}
\usepackage{microtype}
\usepackage{graphicx}
\usepackage{booktabs}
\usepackage[hidelinks]{hyperref}
\usepackage{titlesec}
\titlespacing*{\section}{0pt}{1.2ex}{0.6ex}

\title{\vspace{-2em}Your Project Title}
\author{%
  Group name: \textit{Your Group Name}\\[0.4em]
  First Last (uni1) \quad First Last (uni2) \quad First Last (uni3) \quad First Last (uni4)}
\date{COMS 6998 Networks for AI \quad Project Proposal \quad Fall 2026}

\begin{document}
\maketitle

\section{Problem and Motivation}
What question or problem are you tackling, and why does it matter for serving
LLMs at scale? Connect it to at least one course topic (e.g., KV cache,
disaggregation, speculative decoding, split or edge inference, VLMs).

\section{Related Work}
Two to four papers or systems your project builds on or differs from; course
readings count. For example, DistServe~\cite{distserve} separates prefill and
decode onto different GPUs.

\section{Approach}
What will you build, measure, model, or analyze? For an engineering project,
describe the system and its main components. For a research project, state
the hypothesis or design and how you will test it.

\section{Evaluation Plan}
What does success look like? List your metrics (e.g., latency, throughput,
bandwidth, memory, accuracy), your baseline, and the workloads or models you
will use.

% Optional: a table of planned experiments.
% \begin{tabular}{lll}
%   \toprule
%   Experiment & Metric & Baseline \\
%   \midrule
%   ... & ... & ... \\
%   \bottomrule
% \end{tabular}

\section{Resources and Risks}
Models, datasets, and compute you need (e.g., the course T4 GPU, CPU only,
your own hardware). What is the biggest risk, and what is your fallback if it
goes wrong?

\section{Plan and Roles}
Rough milestones up to the presentations, and who does what.

\begin{itemize}
  \item \textbf{Week of Oct 26:} ...
  \item \textbf{Week of Nov 2:} ...
  \item \textbf{Presentations:} ...
\end{itemize}

% References: add entries below and cite them with \cite{key}.
\begin{thebibliography}{9}
\small
\bibitem{distserve}
Y.~Zhong et al.
\newblock DistServe: Disaggregating prefill and decoding for goodput-optimized
  large language model serving.
\newblock In \emph{OSDI}, 2024.
\end{thebibliography}

\end{document}
